%================================
% Vocabulary & Language Focus — XB1 (S–T)
%================================
\documentclass[../main.tex]{subfiles}

\begin{document}

%================================
\section{S}
%================================

\vw{sail}{/seɪl/}{vi.\&vt.\ \&\ n.}

\sense{vi. \& vt.}{（乘船）航行；起航；驾驶（船只）}
\col{sail from / to / across}
\ex{The fleet sailed from Portsmouth at dawn.}
\col{sail into / out of harbor}
\ex{We watched the yachts sail gracefully out of the harbor.}
\col{go sailing（去驾船游玩）}
\ex{They go sailing on the lake most summer weekends.}

\sense{vi.}{飘，掠}
\ex{Clouds sailed slowly across the autumn sky.}
\ex{A ball of paper sailed across the classroom and landed on the teacher's desk.}

\sense{n.}{帆；乘船航行}
\col{under sail（扬帆航行中）}
\ex{The replica ship crossed the Atlantic under sail alone.}
\col{set sail (for)（启航驶往）}
\ex{The ship set sail for Shanghai before daybreak.}
\col{a sail boat / a three-hour sail（帆船／三小时航程）}
\ex{It's a two-hour sail to the island.}

\begin{gram}{sail / sail for / set sail}
三个高频搭配勿混：\textit{sail \textbf{from} Qingdao \textbf{to} Dalian}（从……航行至……）；\textit{set sail \textbf{for} Dalian}（启航\textbf{驶往}——用 for 不用 to，与 leave for / head for 同一“目的介词”逻辑）；\textit{under sail}（状态）。派生词 \textbf{sailor}（水手）与复合词 \textbf{sailing / sailboat / sails}（帆船运动／帆船／船帆）。俚语 \textit{sail through sth}“轻松通过”（=\,breeze through，见 breeze 词条），考试与体检语境常用。
\end{gram}

\begin{origin}
古英语 \textit{segel}“帆”（原始日耳曼语 *\textit{seglą}），词源可追至原始印欧语 *\textit{sek-}“\textbf{切}”——帆本是“\textbf{从整幅布上裁下的一片}”（同根 \textbf{saw}（锯$\to$切削之器）、拉丁 \textit{secare}“切”$\to$\textbf{section}（切片$\to$章节）、\textbf{sector}、\textbf{scissors}（剪切者）、\textbf{insect}（切入者$\to$昆虫“身体分节”）、\textbf{dissect}（切开$\to$解剖））。由“帆”引申出“航行”，正如汉语“扬帆”“一帆风顺”。看 sail 与 scissors 竟是亲戚，便知“裁一片布”如何撑起了整个大航海时代。
\end{origin}

\vw{scale}{/skeɪl/}{n.}

\sense{n.}{规模，范围；等级，级别；标度；比例}
\col{on a large / small / massive scale（以……的规模，高频）}
\ex{The disaster relief operation was mounted on a massive scale.}
\col{a scale of 1 to 10 / pay scale（十级量表／工资级别）}
\ex{On a scale of 1 to 10, how nervous are you?}
\col{the Richter scale（里氏震级）；a scale model / map（按比例缩小的模型／地图）}
\ex{The earthquake measured 7.8 on the Richter scale.}
\ex{The museum keeps a 1:50 scale model of the ancient capital.}
\col{scale up / down（动词：按比例放大／缩减）}
\ex{The factory had to scale down production during the shortage.}

\begin{gram}{scale“一形三源”：英语里的同形异源词}
\textbf{scale} 是三个来源不同、意义各异的词挤在同一拼写里：（1）本条的“梯级、比例”$\leftarrow$ 拉丁语 \textit{scalae}“\textbf{梯子}”（\textit{scandere}“攀爬”的名词：同根 \textbf{scan}（攀读$\to$扫描）、\textbf{ascend / descend}（上／下爬）、\textbf{transcend}（爬越$\to$超越））——刻度如梯级，“规模”即梯子跨得多大；（2）“鱼鳞”$\leftarrow$ 古法语 \textit{escale}“壳、鳞”（日耳曼源 *\textit{skaljō}“壳”，与 \textbf{shell} 同根）；（3）“秤盘”（英式 weigh 秤）$\leftarrow$ 古诺斯语 \textit{skál}“碗、杯”（与 \textbf{skull} 头骨、\textbf{skoal} 干杯同源）。阅读遇 scale 须先问：梯子、鱼鳞，还是碗？语境定源。
\end{gram}

\vw{scenery}{/ˈsiːnəri/}{n.}

\sense{n.}{风景，景色}
\col{beautiful / spectacular / mountain scenery}
\ex{The train ride offers spectacular mountain scenery.}
\col{scenery（不可数！）：a piece of scenery}
\ex{Rice terraces are the most familiar piece of scenery in his paintings.}

\sense{n.}{舞台布景}
\ex{The scenery was painted to look like a Paris street.}

\begin{lex}[scenery / landscape / view / scene / sight]
\textbf{scenery} 不可数，指一地\textbf{整体自然风貌}（且须自然景，城市风貌不叫 scenery）；\textbf{landscape} 可数，指“大地景致”的整体画面，兼有绘画术语义与喻义“格局”（见 L 部分）；\textbf{view} 是“从某处\textbf{看到}的景”（\textit{a room with a view}）；\textbf{scene} 是具体一“幕”场景，兼指（事件的）现场（\textit{the scene of the accident}）；\textbf{sights} 复数指\textbf{名胜}（\textit{see the sights} 观光）。一句串记：从观景台看到 view，饱览 scenery，拍摄 landscape，路过事故 scene，游览 sights。
\end{lex}

\begin{origin}
经由拉丁语 \textit{scena} $\leftarrow$ 希腊语 \textit{skēnē}“\textbf{帐篷、戏棚}”——古希腊剧场起初不过是\textbf{换装用的帐篷}，后来帐篷前成为表演区，帐篷的墙面即最早的“布景”，遂有 \textbf{scene}（场景）、\textbf{scenery}（布景$\to$自然景，18 世纪“如画风景”审美观的转喻）、\textbf{scenic}（风景优美的：\textit{a scenic route}）、\textbf{scenario}（意大利语借入“情节梗概”$\to$假想情形：\textit{worst-case scenario} 最坏情形）。汉语“风\textbf{景}”本义“日光”，英语“风景”本义“戏棚”——两种语言对“看什么”的原始想象，一个向天，一个向台。
\end{origin}

\vw{security}{/sɪˈkjʊərəti/}{n.}

\sense{n.}{安全；保护措施；保安部门；担保，保证}
\col{national / public security（国家安全／公共安全）}
\ex{The new law strengthens national security and data protection.}
\col{security measures / checks / guards（安保措施／安检／保安）}
\ex{All visitors must pass strict security checks.}
\col{a sense of security（安全感，心理高频）}
\ex{A steady job gave her a sense of security for the first time.}
\col{social security（社会保障）／ security deposit（押金）}
\ex{Social security payments rise with inflation.}
\col{secure（形容词：安全的；动词：获得；固定）}
\ex{She secured a place at a top university.（动词：获得）}

\begin{origin}
拉丁语 \textit{securus}“\textbf{无忧的}”：\textit{se-}（\textit{sine}“\textbf{无}”）$+$ \textit{cura}“\textbf{忧心}”——字面“\textbf{没有牵挂}”（sine cura）。安全的最高境界在罗马人看来不是高墙，而是\textbf{心无挂碍}。\textit{cura}（原始印欧语 *\textit{kwē-}“留心、关切”）家族兴旺：\textbf{cure}（治好“消除那份挂心”）、\textbf{curious}（满心关切$\to$好奇的）、\textbf{curiosity}（好奇心）、\textbf{accurate}（\textit{ad-}+cura 用心照料$\to$精确的）、\textbf{manicure / pedicure}（手／足之护）、\textbf{procure}（代为操办$\to$取得）。“安全”与“好奇”“精确”同根，皆因它们都是\textbf{心的不同用法}。
\end{origin}

\vw{setting}{/ˈsetɪŋ/}{n.}

\sense{n.}{环境，背景；情节背景}
\col{the setting of a novel / film（故事背景）}
\ex{The novel's setting is a small town during the war.}
\col{an ideal / peaceful setting for（……的理想环境）}
\ex{The lake provides an ideal setting for the summer camp.}
\col{a beach / mountain / rural setting}
\ex{The wedding took place in a lovely garden setting.}

\begin{gram}{set 的短语家族与 setting 双义}
\textbf{setting} 既指\textbf{故事背景}（place and time），又指机械\textbf{挡位／设定值}（\textit{change the oven settings} 调烤箱档位；\textit{privacy settings} 隐私设置——数字时代高频）。动词 set 家族庞大，本讲义已见 \textbf{set aside}（见 aside 词条）、\textbf{set sb apart}（见 apart 词条），再加 \textbf{set out / set off}（出发；着手）、\textbf{set up}（建立）、\textbf{set an example}（树立榜样）、\textbf{set the table}（摆放餐具——此义即 setting 的“一套餐具”古义的来源）。凡 set 短语，先辨字面“放置”与喻义“确立”。
\end{gram}

\vw{shade}{/ʃeɪd/}{n.}

\sense{n.}{阴影部分；色度；阴凉处}
\col{in / under the shade of（在……的树荫下）}
\ex{The travelers rested in the shade of a huge camphor tree.}
\col{a shade（一点，些微——口语量词！）}
\ex{The results were a shade disappointing.}
\col{shades of meaning（意味的细微差别）；a shady spot（阴凉处）}
\ex{Dictionaries struggle to capture the subtle shades of meaning.}

\begin{gram}{shade / shadow：一对分分合合的兄弟}
古英语 \textit{scadu} 一词两形：\textbf{shade} 偏\textbf{无形的荫}（阴凉、色调、细微差别——遮阳的功能面），\textbf{shadow} 偏\textbf{有形的影}（人影、影子——物体的轮廓投射）。\textit{sit in the shade}（乘凉，用 shade）；\textit{my shadow grew longer}（影子，用 shadow）。喻义亦分道：\textit{shades of gray}（灰色地带）对 \textit{a shadow of doubt}（一丝怀疑）、\textit{beyond a shadow of a doubt}（毋庸置疑）。而拉丁语的“影”更妙：\textit{umbra}$\to$\textbf{umbrella}（小雨伞“荫具”）、\textbf{umbrage}（荫$\to$不快：\textit{take umbrage} 生气）、\textbf{penumbra}（半影，天文术语）——罗马人连雨伞都是“影子”做的。
\end{gram}

\vw{sorrow}{/ˈsɒrəʊ/}{n.}

\sense{n.}{悲伤，悲痛；伤心事}
\col{express / share / drown one's sorrow（表达／分忧／借酒消愁）}
\ex{He went to drown his sorrows after losing the final.}
\col{in sorrow（悲伤地）；to sb's sorrow（令某人痛心的是）}
\ex{To her great sorrow, the old theater was pulled down.}
\col{sorrowful / joys and sorrows（悲伤的／悲欢）}
\ex{Life is woven of joys and sorrows.}
\ex{She threw a sorrowful glance at the empty chair.}

\begin{origin}
古英语 \textit{sorg}“\textbf{忧虑、悲伤}”（原始日耳曼语 *\textit{surgō}，与德语 \textbf{Sorge}“忧虑”同根——海德格尔哲学关键词 \textit{Sorge}“操心／牵挂”即此词！弗洛伊德的 \textit{Sorge} 一系亦通）。词源或出原始印欧语 *\textit{swergh-}“操心、患病”。英语 \textbf{care}（在乎；忧虑）走过平行路径：由“牵挂”而“照料”而“忧愁”（\textit{careworn} 忧心忡忡的、\textit{for all I care}、\textit{who cares}）。汉语“愁”字从心秋声——草木摇落的季节压在心上；德英的“悲”则长在\textbf{操心}上。三条词源路，都通向心。
\end{origin}

\vw{sour}{/ˈsaʊər/}{adj.\ \&\ vi.\&vt.}

\sense{adj.}{酸的，有酸味的；馊的}
\col{sour milk / yogurt（酸；馊）}
\ex{The milk has gone sour --- throw it away.}
\col{a sour taste in the mouth（嘴里酸涩感；喻“不快余味”）}
\ex{The argument left a sour taste in the mouth.}

\sense{vi. \& vt.}{（使）变坏；变味，酸腐}
\col{go sour / sour（关系、局势）变糟}
\ex{Relations between the two families soured over the land dispute.}
\col{sour sb on sth（使某人对……反感，美式）}
\ex{The long delays soured investors on the project.}

\begin{gram}{sour grapes 与酸味隐喻}
伊索寓言《狐狸与葡萄》留给英语 \textbf{sour grapes}“酸葡萄心理”（贬低得不到之物以求自慰）：\textit{He says the prize doesn't matter, but that's just sour grapes.} 味觉通感（见 bitter 词条）在 sour 上格外多产：\textit{a sour mood / a sour face}（坏情绪／苦脸）、\textit{sour note}（刺耳走调的音）。反差幽默的 \textit{sweet and sour}（糖醋）则是中餐西渐的活词汇。写作中“态度转坏”首选动词 sour，比 become negative 简活。
\end{gram}

\vp{speak volumes about}

\sense{v. 短语}{充分说明，清楚表明}
\ex{Her calm under pressure speaks volumes about her character.}
\ex{The empty shelves spoke volumes about the state of the economy.}

\begin{gram}{volume 的“大”与 speak volumes}
\textbf{volume} 兼指“音量”“体积”“量、额”与“（书之）卷”——四义同源（见下框），故 \textit{speak volumes} 字面“说出好几个卷帙”即“所言极多、信息量极大”，与中文“\textbf{无声胜有声}／不言自明”异曲同工。注意 volumes 用\textbf{复数}；同义升级链：\textit{show $\to$ indicate $\to$ suggest $\to$ \textbf{speak volumes about}}，语气递强，议论文收束句利器：\textit{These figures speak volumes about the cost of the war.}
\end{gram}

\begin{origin}
经由法语 \textit{volume} $\leftarrow$ 拉丁语 \textit{volumen}“\textbf{卷起之物}”＝\textit{volvere}“\textbf{滚动}”（原始印欧语 *\textit{wel-}“转、卷”）的中世纪被动分词——古书是\textbf{卷起的莎草纸}，一“卷”即一书。\textit{volvere} 的滚轮家族在英语中全速转动：\textbf{revolve}（转回去$\to$旋转）、\textbf{evolve}（转出来$\to$演化）、\textbf{involve}（卷进去$\to$涉及）、\textbf{revolt}（卷身翻起$\to$反抗）、\textbf{vault}（拱顶“翻转之穹”）、\textbf{wallow}（打滚，日耳曼表亲）、\textbf{wallet}（小卷包$\to$钱包）、\textbf{volume} 的动词兄弟 \textbf{convoluted}（卷在一起的$\to$盘根错节的）。“音量”义出自书本之“份量”$\to$“量”之义链。
\end{origin}

\vw{stability}{/stəˈbɪləti/}{n.}

\sense{n.}{稳定（性），稳固（性）}
\col{political / economic / social stability}
\ex{Investors fled the country in search of political stability.}
\col{give / provide stability（给予稳定）}
\ex{The fixed pension gives the family some financial stability.}
\col{stable / stabilize（形容词／动词：使稳定）}
\ex{The patient's condition is stable but critical.}
\ex{The bank raised interest rates to stabilize the currency.}

\begin{gram}{stā-“站”词族（三访）：站稳即稳定}
\textbf{stability} $\leftarrow$ 拉丁语 \textit{stabilis}“站得住的”（\textit{stare}“站”，词族详见 consist 词条，本讲义已三访此根）：站得住 = 稳。家族补遗：\textbf{staple}（站柱$\to$主要产品；订书钉“固定之物”）、\textbf{stabilizer}（稳定器；船的减摇鳍）、\textbf{establish}（e- 立起来$\to$建立）、\textbf{obstinate}（站着不走$\to$顽固的）。心理学 \textit{emotional stability}（情绪稳定性）与生态学的 \textit{ecosystem stability}（生态系统稳定性）是科普语篇双料高频。
\end{gram}

\vw{stare}{/steər/}{vi.\ \&\ n.}

\sense{vi. \& n.}{盯着看，凝视，注视}
\col{stare at sb/sth（盯着）}
\ex{The child stared at the magician in open-mouthed wonder.}
\col{stare into the distance / space（出神远望）}
\ex{She stared into space, lost in memories.}
\col{give sb a blank / cold stare}
\ex{He answered the question with a cold stare.}
\col{stare sb in the face（就在眼前；显而易见）}
\ex{The answer had been staring us in the face all along.}

\begin{gram}{“看”的礼节课：stare / gaze / glance / glimpse / glare}
\textbf{stare}（盯着看，常失礼或因惊愕）；\textbf{gaze}（凝望，长久而柔和：\textit{gaze at the stars}）；\textbf{glance}（匆匆一瞥：\textit{glance at one's watch}）；\textbf{glimpse}（瞥见，结果性：\textit{catch a glimpse of}）；\textbf{glare}（怒视：\textit{glare at}）；\textbf{witness}（目击，事件性）。完形填空的经典考区：他\textbf{瞥了一眼}表（glance），她\textbf{怒目}而视（glare），人群\textbf{瞪}着魔术师（stare）——看的方式即态度，选词即选情。
\end{gram}

\vw{stimulate}{/ˈstɪmjuleɪt/}{vt.}

\sense{vt.}{促进，激发；刺激，使兴奋}
\col{stimulate growth / demand / the economy（促进）}
\ex{The tax cuts were designed to stimulate investment.}
\col{stimulate interest / curiosity / discussion（激发）}
\ex{The debate stimulated fresh interest in local history.}
\col{stimulating（令人振奋的；启发性的）}
\ex{It was a stimulating conversation that went on past midnight.}

\begin{gram}{stimulate / encourage / spur / boost}
“促进”四重奏：\textbf{encourage} 最宽（给人以勇气：\textit{encourage sb to do}）；\textbf{stimulate} 强\textbf{激活内在反应}（生理或思维：刺激神经、激发讨论）；\textbf{spur}（马刺“扎一下”$\to$鞭策：\textit{spur sb on to greater efforts}）；\textbf{boost}（从下一托$\to$提升，配数字尤佳：\textit{boost sales by 20\%}）。写作搭配锁：interest/curiosity 配 stimulate，confidence 配 boost，sb 配 encourage/spur。
\end{gram}

\begin{origin}
拉丁语 \textit{stimulare}“\textbf{以尖物驱赶}”，名词 \textit{stimulus}“\textbf{赶牲口的尖棍}”——经济学与心理学的术语 \textbf{stimulus}（刺激；刺激计划 \textit{stimulus package}）、复数 \textbf{stimuli}（遵循拉丁复数，如 data、criteria），词源竟是一根\textbf{戳牛棍}。同根 \textbf{stimulant}（兴奋剂，体育新闻常客：\textit{banned stimulants} 禁用兴奋剂）、\textbf{goad}（赶牛刺棒，日耳曼对词，喻“驱动力”）。从牛棚到美联储的救市方案，一根棍子走了两千年。
\end{origin}

\vw{stream}{/striːm/}{vi.\&vt.\ \&\ n.}

\sense{vi. \& vt.}{鱼贯而行，涌动；流，流动}
\col{stream into / out of（人流涌动）}
\ex{Fans streamed out of the stadium singing and dancing.}
\ex{Tears streamed down her cheeks.（泪如泉涌）}
\col{be streaming with（雨水、汗水淌下）}
\ex{The windows were streaming with rain.}

\sense{n.}{小溪；（人）流，（车）流}
\col{a stream of sth（一连串的）}
\ex{A steady stream of visitors filed past the coffin.}
\ex{He kept up a constant stream of jokes during the trip.}
\col{go / swim against the stream（逆流而行，喻不从众）}
\ex{In that decade, choosing art was swimming against the stream.}

\begin{gram}{stream 的新义：流媒体时代}
数字时代给 stream 插上了翅膀：\textbf{stream sth}“\textbf{在线播放／收听}”（\textit{stream a concert live} 直播演唱会；\textit{streaming services} 流媒体平台），名词 \textbf{mainstream}（主流：\textit{mainstream media / values}）与其反义 \textbf{off-stream?}（否——惯用 \textbf{alternative / fringe}）构成文化讨论的常用轴。一词从山涧流进了互联网，是“旧词新义”速度的标本。
\end{gram}

\vw{strike}{/straɪk/}{vi.\&vt.\ (struck, struck)\ \&\ n.}

\sense{vt.}{突然想到；给……以印象}
\col{it strikes sb that\dots（某人突然想到，无灵主语句式）}
\ex{It suddenly struck her that she had seen the man before.}
\col{strike sb as (being) adj.（给某人……的印象）}
\ex{The answer struck him as too simple to be true.}

\sense{vt. \& vi.}{撞，碰；打；划（火柴）；罢工；袭击}
\col{strike a match（划火柴）／ strike the hour（钟鸣报时）}
\ex{The grandfather clock struck twelve.}
\col{go on strike / strike for better pay}
\ex{Nurses went on strike for a pay rise.}
\col{be struck by lightning（遭雷击）}
\ex{The old oak was struck by lightning in the storm.}
\col{strike a blow for / against（为／反对……而斗争）}
\ex{The verdict struck a blow for press freedom.}

\sense{n.}{罢工；袭击；击，打}
\col{an air strike（空袭）；a general strike（总罢工）}
\ex{Air strikes targeted the enemy's supply lines.}

\begin{gram}{strike 的“无灵主语”与 strike a balance}
两条高级句式：（1）\textit{\textbf{It struck me that\dots}}（我忽然想到）——无灵主语 it 承载思维活动，比 I suddenly thought 更利落，同族 \textit{it occurs to me that}；（2）\textbf{strike a balance between A and B}（在 A 与 B 之间求得平衡）——议论文万用结论句：\textit{We must strike a balance between growth and conservation.} 顺带分清：strike（打、击）$+$ stroke（轻抚；中风；笔画：\textit{breaststroke} 蛙泳）是同根 doublet，一重一轻。
\end{gram}

\begin{origin}
古英语 \textit{strīcan}“\textbf{擦过、击打、行走}”（原始日耳曼语 *\textit{strikaną}，与德语 \textit{streichen}“抚、涂”同根；或出原始印欧语 *\textit{sterig-}“僵、击”）。它的一次词义大爆炸覆盖英语半部日常语：\textbf{strike}（击）之外，\textbf{stroke}（轻击$\to$抚摩；中风“脑被击”）、\textbf{streak}（擦痕$\to$条纹；\textit{a streak of luck}）、\textbf{trace / track}?（另一词源）、\textbf{trick}?（另源）。最妙的是\textbf{罢工}义：1768 年伦敦船工\textbf{strike（撤走）}帆桁（strike the topgallant sails“降帆罢工”），水手语的“降下”遂成劳动运动的旗帜——一次航海操作变成了社会史名词。
\end{origin}

\vw{striking}{/ˈstraɪkɪŋ/}{adj.}

\sense{adj.}{引人注目的，显著的；妩媚动人的，标致的}
\col{a striking resemblance / contrast / feature}
\ex{There is a striking resemblance between the two portraits.}
\ex{The most striking feature of the report is what it leaves out.}
\col{a striking woman（引人注目的女子）}
\ex{In her red coat she was a striking figure at the party.}
\col{strikingly（显著地）}
\ex{The two brothers are strikingly different in temperament.}

\begin{gram}{“显著”形容词群：写作的强调度盘}
从弱到强：\textbf{noticeable}（可察觉）$\to$ \textbf{marked / evident}（明显的）$\to$ \textbf{striking}（令人瞩目）$\to$ \textbf{dramatic / remarkable}（戏剧性／非凡）。图表作文的升降描述亦分档：\textit{a slight rise $\to$ a marked rise $\to$ a striking / dramatic rise}。\textbf{strike} 的现在分词形容词化（striking）与过去分词（struck“被震住的”：\textit{dumb-struck} 目瞪口呆）一主动一被动，分工如 interesting/interested。
\end{gram}

\vw{structure}{/ˈstrʌktʃər/}{n.\ \&\ vt.}

\sense{n.}{建筑物；结构，构造；精心组织}
\col{the structure of sth（……的结构）}
\ex{The structure of the novel mirrors the rise and fall of the family.}
\col{a steel / wooden structure（建筑实体）}
\ex{The pavilion is a light wooden structure open on all sides.}
\col{grammatical structure / sentence structure（语法结构）}
\ex{Compare the sentence structure of the two translations.}

\sense{vt.}{精心组织，安排}
\ex{Structure your essay so that each paragraph makes one point.}
\ex{The training program is structured around real projects.}

\begin{gram}{struct-“建造”词族（二访）与“结构化”写作术语}
拉丁语 \textit{struere}“堆积建造”（词族首访见 industrialization 词条）家族补全：\textbf{construct}（合建$\to$建造；构念）、\textbf{destruction}（拆建$\to$毁灭）、\textbf{instruct}（往里筑$\to$教导；指示）、\textbf{instrument}（筑造之具$\to$工具；乐器）、\textbf{obstruct}（ob- 挡路堆$\to$阻碍：\textit{obstruct the view}）。写作课术语 \textbf{structure} 动词化（“谋篇”）是学术英语的地道用法；相关词 \textbf{infrastructure}（infra- 在下$\to$基础设施）、\textbf{restructure}（重组）为新闻高频。
\end{gram}

\vw{suitable}{/ˈsuːtəbl; ˈsjuːtəbl/}{adj.}

\sense{adj.}{合适的，适宜的，适当的，适用的}
\col{be suitable for sb/sth（固定介词 for）}
\ex{This movie is not suitable for young children.}
\col{be suitable to do sth}
\ex{The empty warehouse is suitable to host the exhibition.}
\col{highly / ideally / perfectly suitable}
\ex{The hotel is ideally suitable for families with small children.}

\begin{gram}{suit 的“跟随”身世与 sequi 词族}
\textbf{suitable} $<$ \textbf{suit}“合宜”$<$ 拉丁语 \textit{sequi}“\textbf{跟随}”（原始印欧语 *\textit{sekʷ-}“跟随”，见 association 词条）——\textbf{suit} 的诸义全是“跟随”的变奏：\textbf{一套衣服}（相随的一组）；\textbf{诉讼}（追随对方讨说法：\textit{file a suit against}）；\textbf{suit sb}（“随合”某人$\to$适合）；\textbf{pursue / pursuit}（pur- 向前跟随$\to$追求：\textit{in pursuit of happiness}）；\textbf{suitcase}（随行之箱）；\textbf{suite}（随员；一套房间；组曲——读音 /swiːt/，勿与 suit 混）；\textbf{ensue}（随之而来）。辨析回看 appropriate 词条：suitable 泛指合用，appropriate 重得体。
\end{gram}

\vw{surroundings}{/səˈraʊndɪŋz/}{n.}

\sense{n.}{周围，环境}
\col{pleasant / unfamiliar / new surroundings（恒用复数！）}
\ex{The children soon adapted to their new surroundings.}
\col{in comfortable / safe surroundings}
\ex{The patients recover faster in pleasant surroundings.}
\col{surround（动词：环绕）／ surrounding（形容词：周围的）}
\ex{The school is surrounded by a high wall.}
\ex{The surrounding villages were cut off by the flood.}

\begin{gram}{恒用复数的“境况”名词}
\textbf{surroundings} 无单数用法，同类“环境／境况”词多以复数行世：\textbf{belongings}（所有物）、\textbf{earnings}（收入）、\textbf{savings}（积蓄）、\textbf{findings}（调查发现）、\textbf{wages}（工资）、\textbf{circumstances}（境况：\textit{in difficult circumstances}）、\textbf{conditions}（条件；环境）。谓语动词一律复数。近义辨析：surroundings（物质环境，身边）、environment（生态／社会总环境）、circumstances（人事境遇）、settings（事件背景，见 setting 词条）。
\end{gram}

\begin{origin}
\textbf{surround} $<$ 古法语 \textit{suronder}“\textbf{漫溢}”$<$ 晚期拉丁语 \textit{superundare}：\textit{super-}（上面）$+$ \textit{undāre}“\textbf{起浪}”（名词 \textit{unda}“\textbf{波浪}”，原始印欧语 *\textit{wed-}“湿、水”）。“包围”的本义是\textbf{潮水漫过头顶}。\textit{unda} 的波纹荡满英语：\textbf{abundant}（ab- 离岸起浪$\to$\textbf{丰富}——浪涌而出即充裕）、\textbf{inundate}（in- 涌入$\to$\textbf{淹没；使应接不暇}：\textit{inundated with emails}）、\textbf{undulate}（起伏如波）、\textbf{redound}（回涌$\to$有助于）。而 \textit{unda} 与 \textbf{water}、\textbf{wet} 同出印欧水根 *\textit{wed-}——被包围，就是被水环抱。
\end{origin}

\vw{sweep}{/swiːp/}{vi.\&vt.\ (swept, swept)\ \&\ n.}

\sense{vi. \& vt.}{席卷，横扫；打扫；吹走；清除}
\col{sweep across / through（风潮席卷）}
\ex{A wave of strikes swept across the country that winter.}
\col{sweep the floor / the board（扫地／大获全胜）}
\ex{The young team swept the board, winning every medal.}
\col{be swept away / sweep sb off their feet（冲走／令人倾倒）}
\ex{The film swept audiences off their feet worldwide.}
\col{sweep sth under the carpet（把问题扫进地毯下，喻掩盖）}
\ex{The committee simply swept the complaint under the carpet.}

\sense{n.}{打扫，清扫；挥动；搜查}
\col{give the room a sweep（快速打扫一遍）}
\ex{Give the porch a quick sweep before the guests arrive.}

\begin{gram}{ept 型不规则动词与“风一般”的意象动词}
\textbf{sweep--swept--swept} 与 \textbf{sleep--slept, keep--kept, weep--wept, feel--felt, creep--crept} 同属“长元音 ee $\to$ 短元音 e + t”的强变化类——古英语七类强动词的现代遗存。sweep 的习语几乎都带\textbf{大开大合}的动势：\textit{sweep to power}（以压倒优势上台）、\textit{a sweeping statement}（一竿子打翻一船人的笼统论断）、\textit{sweeping changes}（大刀阔斧的变革）。写作中表“席卷”，sweep 比 spread quickly 更具画面速度感。
\end{gram}

%================================
\section{T}
%================================

\vp{tear apart}

\sense{v. 短语}{使（关系密切的人）分离，分开；撕碎}
\ex{The war tore thousands of families apart.}
\ex{She tore the envelope apart, looking for the letter.}
\col{be torn between A and B（在 A 与 B 之间左右为难，被动高频）}
\ex{He was torn between duty and love.}
\col{be torn apart by（被……撕裂、摧毁）}
\ex{The community was torn apart by the scandal.}

\begin{gram}{同形异音词（heteronym）三连：tear / tear}
\textbf{tear}\textsubscript{1} /tɪər/ $v.$ 撕裂（tear--tore--torn）；\textbf{tear}\textsubscript{2} /teər/ $n.$ 眼泪（恒用复数 \textit{tears}；动词“含泪” \textit{tear up} /ˈtɪər ʌp/?——注意 tear up“撕毁”读 /tɪər/，而“落泪”说 \textit{burst into tears / his eyes teared up}）。同类heteronym 还有 \textbf{wind}（风 /wɪnd/ vs.\ 上发条 /waɪnd/，见 W 部分）、\textbf{close}（关闭 /kləʊz/ vs.\ 亲近的 /kləʊs/）、\textbf{live}（居住 /lɪv/ vs.\ 现场直播的 /laɪv/）。听辨与朗读的暗礁，务必逐对记忆。
\end{gram}

\vp{throw in}

\sense{v. 短语}{奉送，额外赠送}
\ex{If you buy two shirts, the shop will throw in a free tie.}
\ex{Delivery is thrown in free of charge this month.}
\col{throw in the towel / the sponge（认输，弃权——拳击术语）}
\ex{After three failed attempts, the team refused to throw in the towel.}

\vp{throw oneself into}

\sense{v. 短语}{投身于，热衷于}
\ex{After retirement, he threw himself into volunteer teaching.}
\ex{She threw herself into the rehearsals, often staying late into the night.}

\begin{gram}{throw 短语的身姿隐喻}
throw 家族短语的共性是\textbf{整个身体的投入或抛掷}：\textbf{throw oneself into}（纵身投入）；\textbf{throw sth off}（甩脱：\textit{throw off the cold} 摆脱感冒）；\textbf{throw sth together}（匆匆拼凑：\textit{throw together a meal}）；\textbf{throw light on sth}（照亮、阐明——与 cast light on 同义，见 cast 词条）；\textbf{throw a party}（办派对）；\textbf{be thrown into jail}（投入监狱）。写作中 \textit{throw oneself into} 比 engage in 更有体温，续写描写人物专注状态的首选。
\end{gram}

\vw{tolerant}{/ˈtɒlərənt/}{adj.}

\sense{adj.}{宽容的，容忍的；能耐……的}
\col{be tolerant of sb/sth（固定介词 of）}
\ex{A good teacher is tolerant of students' mistakes.}
\ex{Society has become more tolerant of different lifestyles.}
\col{tolerate / tolerance / tolerable（词族）}
\ex{The school does not tolerate bullying in any form.（容忍，多否定）}
\ex{Patients gradually build up a tolerance to the medicine.（耐受性）}
\ex{The noise was tolerable at first, then unbearable.（尚可忍受的）}

\begin{gram}{tolerable / tolerant：-able 与 -ant 的一念之差}
\textbf{tolerant}（主动“能宽容他人的”，形容人）；\textbf{tolerable}（被动“\textbf{可被容忍的}”，形容事物：\textit{a tolerable error}）。同型对子：\textbf{acceptable}（可接受的）对 \textbf{accepting}（包容接纳的）；\textbf{bearable}（可忍受的）对 \textbf{bearing}。联想英语警句：\textit{The tolerant are few; the tolerable are many.} 宽容者少，可容忍者多——一词之差，主客易位，恰与 deserve doing 的主动表被动（见 deserve 词条）互为镜像。
\end{gram}

\begin{origin}
拉丁语 \textit{tolerāre}“\textbf{背负、承受}”（词源或与 \textit{tollere}“举起”相关，原始印欧语 *\textit{telə-}“承载、称量”）。其家族多与“扛”有关：\textbf{tolerate}（扛得住$\to$容忍）、\textbf{tolerance}（承受力；\textbf{tolerance} 在工程学中是“\textbf{公差}”——机器扛得住的误差余量！）、\textbf{tolerant}、\textbf{extol}?（另源 \textit{tollere} 举起$\to$颂扬）。同根的日耳曼远亲是 \textbf{toler}？否——但 \textbf{talent}（才能）出自希腊\textbf{塔兰同}（重量单位$\to$《马太福音》按才受托的比喻$\to$才能）：\textbf{tolerate 与 talent 同出“称量的秤”}，一个扛重量，一个是重量。宗教改革时期的“宽容令” \textit{Edict of Toleration} 使 tolerance 进入政治语汇。
\end{origin}

\vw{transform}{/trænsˈfɔːm/}{vi. \& vt.}

\sense{vi. \& vt.}{（使）变形，（使）改观，（使）转化}
\col{transform sth (from A) into B（固定搭配 into）}
\ex{The new bridge transformed the island from a backwater into a boom town.}
\ex{Three years of training completely transformed her body and mind.}
\col{transform sth beyond recognition（改变得面目全非）}
\ex{The war transformed the city beyond recognition.}
\col{transformation（名词：变革；〔物理〕变换）}
\ex{The digital transformation of the economy is accelerating.}

\begin{gram}{form-“形”词族与 trans- 前缀总动员}
拉丁语 \textit{formāre}“塑形”家族整齐列队：\textbf{transform}（trans- 移形$\to$转化）、\textbf{reform}（re- 回复原形$\to$改革）、\textbf{deform}（de- 毁形$\to$使变形）、\textbf{conform}（con- 共同之形$\to$遵从：\textit{conform to the rules}）、\textbf{inform}（in- 赋形于心$\to$告知；信息）、\textbf{perform}（per- 彻底成形$\to$执行；表演）、\textbf{uniform}（uni- 一形$\to$制服；一致的）。前缀 \textbf{trans-}“越过、转换”另立一军：\textbf{transfer / transport / transmit / translate / transplant / transmit}（移交／运输／传送／翻译／移植）——凡 trans-，必有一场“跨界搬运”。
\end{gram}

\vw{triumph}{/ˈtraɪʌmf/}{n.\ \&\ vi.}

\sense{n.}{巨大成功，重大成就，伟大胜利}
\col{a triumph for / of sth（……的胜利）}
\ex{The vaccine is a triumph of international cooperation.}
\ex{Her return to the stage after the accident was a personal triumph.}
\col{in triumph（得意扬扬地）}
\ex{He carried the trophy off in triumph.}

\sense{vi.}{打败，战胜，成功}
\col{triumph over sb/sth（战胜，书面）}
\ex{In the end, common sense triumphed over prejudice.}

\begin{gram}{triumph / victory / success 的修辞高度}
\textbf{victory} 战胜（军事、竞技的胜利，与 defeat 相对）；\textbf{success} 达成目标（中性、使用最广）；\textbf{triumph} 凯旋式的\textbf{辉煌大胜}（情感与历史感最强，常与 of 连用点出凯旋之物：\textit{a triumph of the human spirit}）。习语 \textbf{triumphal arch}（凯旋门：巴黎 Arc de Triomphe）、形容词 \textbf{triumphant}（得胜的；得意洋洋的：\textit{a triumphant smile}）。写作中把 success 升级为 triumph，一字提升庆典浓度。
\end{gram}

\begin{origin}
经由古法语 \textit{triumphe} $\leftarrow$ 拉丁语 \textit{triumphus}“\textbf{罗马凯旋式}”——将军大胜后，元老院批准他率军\textbf{穿越罗马城}（可能借自希腊语 \textit{thriambos}，酒神狄俄尼索斯的\textbf{凯旋颂歌}，一说是伊特鲁里亚语借词）：战车载着俘获与财宝，阶下囚游街，身后有奴隶低语“记住你是凡人”（\textit{memento mori}）。\textbf{triumph 的词源就是这场游行本身}。同根酒神颂歌还生出 \textbf{thriamb} 与弥尔顿笔下的 \textit{Bacchic triumph}。从酒神的歌到将军的游行再到普通人的“小胜利”，一个词装下了整个罗马的荣耀仪式。
\end{origin}

\vw{tune}{/tjuːn/}{n.\ \&\ vt.}

\sense{n.}{曲调，曲子}
\col{hum / whistle a tune（哼／吹曲子）}
\ex{She hummed a cheerful tune as she cooked.}
\col{in tune / out of tune with（与……合拍／格格不入，喻义高频）}
\ex{His ideas are completely out of tune with the times.}
\col{change your tune（改变论调）}
\ex{When the money ran out, investors quickly changed their tune.}
\col{to the tune of（金额高达，戏谑）}
\ex{The firm lost money to the tune of two million yuan.}

\sense{vt.}{调音；调频道；调整}
\col{tune a piano / an instrument（调音）}
\ex{The pianist was tuning the piano before the concert.}
\col{tune in (to)（收听；关注）}
\ex{Millions tuned in to watch the final episode.}
\col{tune sth up（调校；热身）／ stay tuned（别走开，敬请期待）}
\ex{Stay tuned for the results after the break.}

\begin{gram}{tone / tune：一对分家的孪生词}
\textbf{tune} 本是 \textbf{tone} 的变体（受古法语 \textit{ton} 影响的拼法分化）：tone $\leftarrow$ 拉丁语 \textit{tonus}“\textbf{音调、张力}”$\leftarrow$ 希腊语 \textit{tonos}“\textbf{绷紧}”（\textit{teinein}“绷、拉紧”，原始印欧语 *\textit{ten-}“拉伸”——\textbf{tension, tense, tent, extend, tendon, tendril}、与 \textit{tone} 的兄弟 \textbf{hypertony?} 无关，但 \textbf{baritone}（男中音“重音”）、\textbf{monotone}（单音调）皆是）。分家后分工：\textbf{tone} 管“音色、语气、色调”（\textit{in an angry tone; the warm tones of the painting}），\textbf{tune} 管“旋律、调子”（\textit{a catchy tune}）。正如汉语“调”一字兼 dìao（曲调）与 tóng?（音色称“音色”），英语把一根绷紧的弦分成了两支。
\end{gram}

%================================
\section{U}
%================================

\vw{undergo}{/ˌʌndəˈɡəʊ/}{vt.\ (underwent, undergone)}

\sense{vt.}{经历，经受（变化、手术、考验等重大过程）}
\col{undergo treatment / surgery / an operation}
\ex{He underwent emergency surgery after the crash.}
\col{undergo changes / transformation / tests}
\ex{The company has undergone a complete transformation in five years.}
\ex{All new products undergo rigorous safety tests.}

\begin{gram}{under- 复合动词家族与 undergo 的“沉重”语域}
\textbf{undergo} $=$ under + go“从下面走过”，喻“\textbf{负重而行、承受}”——故其宾语几乎总是\textbf{重大而不轻松的过程}（手术、考验、变革），\ding{55} \textit{undergo a holiday}。同族：\textbf{undertake}（承担；从事：\textit{undertake a task}，undertook--undertaken）、\textbf{underline}（画线$\to$强调）、\textbf{underestimate}（低估——见 estimate 词条）、\textbf{undermine}（挖墙脚$\to$暗中削弱：\textit{undermine confidence}）、\textbf{underlying}（底层的：\textit{the underlying cause} 根本原因）。under- 家族的意象总纲：\textbf{在下、不足、承受}。
\end{gram}

\vw{unconstrained}{/ˌʌnkənˈstreɪnd/}{adj.}

\sense{adj.}{不受约束的，自由的}
\col{unconstrained by sth（不受……束缚，被动式常见）}
\ex{His style, unconstrained by tradition, blended jazz with local opera.}
\col{unconstrained growth / spending（无节制的增长／花销）}
\ex{Unconstrained growth of tourism is damaging the old town.}
\col{free and unconstrained（自由自在，汉语式叠用对应）}
\ex{The children ran free and unconstrained across the grass.}

\begin{gram}{strain“捆紧”词族回访与 constrained 的写作用途}
\textbf{unconstrained} $=$ un-（不）$+$ con-（共同）$+$ \textit{stringere}“\textbf{捆紧}”的过去分词——字面“\textbf{没有被捆住的}”。这个词根已两次到访（见 district 词条）：\textbf{constrain}（捆住$\to$\textbf{限制}：\textit{be constrained by time and money}）、\textbf{constraint}（束缚；限制因素，学术语篇高频）、\textbf{restrain}（re- 拉回捆住$\to$\textbf{克制}：\textit{restrain one's anger}）、\textbf{restraint}（克制；安全带；约束装置）、\textbf{strain}（绷紧；拉伤； strain every nerve）、\textbf{stress}（压力）、\textbf{strict}（捆得紧的$\to$严格的）、\textbf{stringent}（绷紧的$\to$〔规定〕严苛的）。写作对句：\textit{constrained by budget, unconstrained by imagination}——受预算之束，无想象之羁。
\end{gram}

\vw{unforgiving}{/ˌʌnfəˈɡɪvɪŋ/}{adj.}

\sense{adj.}{棘手的；不饶人的，不宽容的}
\col{an unforgiving climate / environment / terrain（严酷无情的自然环境，科考语篇高频）}
\ex{The plateau is an unforgiving environment for any motor vehicle.}
\ex{Antarctica is beautiful but unforgiving.}
\col{an unforgiving schedule / deadline（毫不容情的日程／死线）}
\ex{The athletes trained on an unforgiving schedule before the Games.}
\col{an unforgiving surface（容不得差错的场地、材质）}
\ex{Ice is an unforgiving surface --- one slip means a broken ankle.}

\begin{origin}
\textbf{unforgiving} $=$ un- $+$ \textbf{forgive}“原谅”的现在分词“\textbf{不给予宽恕的}”——自然与人皆可如此。\textbf{forgive} 的构词藏着一段语言接触史：前缀 \textit{for-}（古英语 \textit{for-}“\textbf{彻底、离去}”，同见于 \textbf{forbid}（彻底命令$\to$禁止）、\textbf{forlorn}（被彻底舍弃的$\to$凄凉的）、\textbf{forsake}（彻底放下$\to$抛弃））$+$ \textit{giefan}“给”——“\textbf{彻底地给出去}”。这个组合竟是古英语对\textbf{拉丁语 \textit{perdonare}} 的仿译（\textit{per-}“彻底”$+$\textit{donare}“给予”）：即 \textbf{forgive 与 pardon 是同一隐喻的两种语言实现}——“原谅”在东西罗马与盎格鲁撒克逊人的头脑里，都是\textbf{把怨债整个送出去}。形容词化后，unforgiving 又把这层“给予”反赠给大自然：不饶人的冻原，一视同仁。
\end{origin}

\vw{unusual}{/ʌnˈjuːʒuəl/}{adj.}

\sense{adj.}{特别的，不寻常的；独特的，与众不同的}
\col{an unusual experience / hobby / name}
\ex{Eating insects is an unusual source of protein in some cultures.}
\col{it is unusual for sb to do sth（罕见）}
\ex{It is unusual for this river to freeze even in January.}
\col{unusually（罕见地；异常地）}
\ex{Unusually for him, the great novelist was late.}

\begin{gram}{两个 un-：形容词否定前缀与动词逆转前缀}
英语有\textbf{两个来源不同}的 un-：（1）\textbf{形容词 un-}（日耳曼古语）：unusual, unhappy, unclear, unequal——表“\textbf{非、不}”；（2）\textbf{动词 un-}（稍晚）：untie（解开）、undo（撤销）、unfold（展开）、unveil（揭幕）、unearth（掘出）——表“\textbf{逆动作、还原}”。二者可叠加于同根：\textbf{do}（做）$\to$ \textbf{undo}（撤销）$\to$ \textbf{undone}（未完成的／松开的）。辨异技巧：后接形容词是“不”，后接动词是“反”。相关高频词 \textbf{unveil}（揭去面纱$\to$发布、揭幕，新闻标题常客）。
\end{gram}

\begin{origin}
\textbf{un-} $+$ \textbf{usual}“寻常的”$\leftarrow$ 拉丁语 \textit{ūsualis}，出自 \textit{ūsus}“\textbf{使用、习惯}”（动词 \textit{ūtī}“用”，原始印欧语 *\textit{h₂eu-}? 词源诸说不一，或与“愉悦地享用”相关）。\textit{ūtī} 的家族：\textbf{use}（用）、\textbf{useful}、\textbf{usage}（用法；惯例）、\textbf{usual / usually}、\textbf{unusual}、\textbf{abuse}（ab- 用偏$\to$\textbf{滥用；虐待}）、\textbf{utensil}（用具）、\textbf{utilize}（utilitas$\to$利用，语体最正式）。注意 \textbf{usual} 读 /ˈjuːʒuəl/ 而 \textbf{unusual} 的 -su- 读 /ʒ/ 由 -sual 逆同化而来——“寻常”的反面，连发音也特别。
\end{origin}

\vw{uplifting}{/ˌʌpˈlɪftɪŋ/}{adj.}

\sense{adj.}{令人振奋的，鼓舞人心的}
\col{an uplifting story / speech / experience}
\ex{The documentary ends on an uplifting note.（在振奋的调子上收尾）}
\ex{Her speech was uplifting, full of hope and quiet humor.}
\col{lift sb's spirits（使振作——同义动词短语）}
\ex{The good news lifted everyone's spirits.}

\begin{gram}{“情绪升降”的方位隐喻：uplifting / uplifting vs.\ depressing}
英语以\textbf{上下}喻情绪高低，成对呼应：\textbf{uplifting / uplifting}（扬起$\to$振奋）对 \textbf{depressing}（压下$\to$令人沮丧，见 depress 词条）；\textbf{cheer sb up}（振作）对 \textbf{let sb down}（使失望）；\textbf{in high spirits} 对 \textbf{feel low / down}。而 \textit{lift} 一词（电梯；举起）与 \textit{shoplift}（顺手“举”走$\to$入店行窃）、\textit{lift a passage}（剽窃段落——学术写作警示词）共享“提起”义。跨文化彩蛋：这一“上好下坏”的隐喻为人类各语言所共享（汉语“高涨／低落”），是认知语言学的教科书案例。
\end{gram}

\vw{upset}{/ˌʌpˈset/}{adj.\ \&\ vt.\ (upset, upset)}

\sense{adj.}{难过的，失望的}
\col{be / get upset about / by sth（因……而难过，介词 about/by）}
\ex{She was deeply upset by the criticism.}
\ex{Don't get so upset about a minor setback.}
\col{be upset with sb（生某人的气）}
\ex{He was upset with himself for missing the penalty.}

\sense{vt.}{使烦恼；搅乱；打翻}
\ex{The sudden policy change upset all our plans.（搅乱）}
\ex{The cat upset the vase and fled.（打翻）}
\col{an upset stomach（肠胃不适，名词作定语）}
\ex{Too much ice cream gave him an upset stomach.}

\begin{gram}{set 短语的“反转构词”一族：upset / set up / outset}
\textbf{upset} $=$ set up（安放）的\textbf{语序倒转}，义也随之反转：“掀翻、扰乱”。古英语如此倒装造词的还有：\textbf{outset}（set out 出发的名词“\textbf{开端}”$\leftrightarrow$ \textbf{outset / outset}；与 \textbf{outcome}（come out 的名词“结果”）成对——\textit{from the outset to the outcome} 自始至终）、\textbf{upset} 的名词义“爆冷门”（体育新闻：\textit{a major upset} 大冷门；动词 \textit{upset the champions} 掀翻冠军）。三词皆“短语倒转成名/动词”的化石，辨认时倒着读即可。
\end{gram}

%================================
\section{V}
%================================

\vw{variation}{/ˌveəriˈeɪʃn/}{n.\ \&\ vary /ˈveəri/ vi.\&vt.}

\sense{n.}{变奏，变奏曲；变化，变更；变体，变种}
\col{a variation on sth（……的变奏／变体）}
\ex{The pianist played a brilliant variation on a folk tune.}
\col{regional / seasonal / genetic variation（地区／季节／基因变异）}
\ex{There is wide regional variation in dialect and food.}
\ex{Genetic variation helps species survive disease.}

\sense{vary vi.\&vt.}{（根据情况）变化，变更；相异，不同；变更，改变}
\col{vary from A to B / vary in size / price / quality}
\ex{Opinions on the plan vary from enthusiasm to open hostility.}
\ex{The temperature varies between minus five and plus ten in spring.}
\col{vary with / according to（随……而变，高频）}
\ex{Prices vary with the season and the size of the catch.}
\col{various / variety（词族：各种各样的；种类）}
\ex{a variety of = varieties of（种种）}
\ex{Variety is the spice of life.（谚：变化是生活的调味品）}

\begin{gram}{vary / various / variable / varied 四分}
\textbf{vary}（动词：变化）；\textbf{various}（形容词：\textbf{多个不同的}，修饰可数名词：\textit{various reasons}）；\textbf{varied}（形容词：\textbf{形形色色的、有变化的}：\textit{a varied career / menu}）；\textbf{variable}（形容词：\textbf{易变的}、\textbf{可变的}；名词：\textbf{变量}，数理语篇高频：\textit{an independent variable} 自变量）。写作避复三连：\textit{vary the diet $\to$ a varied menu $\to$ various dishes}，一族的四个面貌各就其位。
\end{gram}

\begin{origin}
拉丁语 \textit{variāre}“使变化”，出自 \textit{varius}“\textbf{斑驳的、杂色的、多变的}”——本义是\textbf{花斑、斑点}（一只羽毛斑驳的鸟就是 \textit{avis varia}）。由此看透全家福：\textbf{vary / variety / variation}（“变化”生于“花色”）、\textbf{variegated}（斑驳的，园艺与生物学术语）、\textbf{variegate}、法语借词 \textbf{varie}——甚至 \textbf{irregular}?（另源）。“变化”的词源是\textbf{色彩斑驳}：先看见羽毛上的花斑，才抽象出“多样”的概念。这与汉语“\textbf{斑斓}”“陆离”以色写杂，竟不谋而合。
\end{origin}

\vw{vast}{/vɑːst/}{adj.}

\sense{adj.}{大量的，辽阔的，巨大的}
\col{a vast area / crowd / majority（辽阔地区／人海／绝大多数）}
\ex{The camel caravans crossed a vast stretch of desert.}
\col{the vast majority of（绝大多数，议论文高频）}
\ex{The vast majority of residents support the new library.}
\col{vast sums / amounts of（巨款／巨量）}
\ex{Vast sums were poured into the space program.}
\col{vastly（大大地）}
\ex{Conditions have improved vastly since the 1990s.}

\begin{origin}
经由拉丁语 \textit{vastus}“\textbf{空旷的、荒凉的}”——\textbf{vast 的本义不是“大”，而是“空”}：人迹全无的旷野之“大”，是无物的可怕之大。它在法语中的兄弟 \textbf{waste}（荒野$\to$\textbf{浪费}；废物；荒废）与 \textbf{devastate}（de- 彻底荒芜$\to$\textbf{摧毁}、名词 \textbf{devastation}）同根（原始印欧语 *\textit{wāhstu-}? 词源与“空”相关）——\textbf{vast 与 waste 是同一个拉丁词在英法两国的双胞胎}，一如 guarantee/warrant。故英语里“浩瀚”与“荒废”血脉相连：\textit{a vast wasteland} 浩瀚荒原，一词族内部自相呼应。
\end{origin}

\vw{version}{/ˈvɜːʃn/}{n.}

\sense{n.}{变体，变种；说法；版本}
\col{the latest version of sth（最新版本）}
\ex{The app updates to the latest version automatically.}
\col{sb's version of events（某人对事件的“说法”，新闻高频）}
\ex{The two drivers gave very different versions of the accident.}
\col{an English / film version（英语版／电影版）}
\ex{The English version of the novel appeared two years later.}
\col{version of the truth（某种“真相”表述，常含贬义）}
\ex{That is his version of the truth, not mine.}

\begin{gram}{verse 词族“转”字诀（vertere 家族）}
\textbf{version} $\leftarrow$ 拉丁语 \textit{vertere}“\textbf{转动}”（原始印欧语 *\textit{wer-}“转、弯”）的过去分词 \textit{versus}“被转成的”——“\textbf{转写而成之物}”：翻译是语言的转写，版本是文本的转向，说法是事实的转向。\textit{vertere} 的旋转家族：\textbf{reverse}（re- 回转$\to$颠倒；倒车）、\textbf{convert}（转过去$\to$转化；改宗）、\textbf{divert}（di- 转开$\to$使改道；转移注意力）、\textbf{diverse}（各自转向$\to$多样的）、\textbf{anniversary}（annus 年 + versus“转回之日”$\to$\textbf{周年纪念}）、\textbf{universe}（uni- 合一 + versus“转成一体之物”$\to$\textbf{宇宙}）、\textbf{versus}（对、对抗，缩写 vs.）、\textbf{verse}（诗行“犁转身回耕的一垄”！）。一个“转”字，从诗句耕到宇宙。
\end{gram}

\vw{vision}{/ˈvɪʒn/}{n.}

\sense{n.}{远见卓识；视力，视野；想象}
\col{a vision of the future（对未来的构想）}
\ex{The company's founding vision was a computer on every desk.}
\col{have a vision of sth（心中看见……）}
\ex{The architect had a clear vision of what the square could become.}
\col{poor / perfect vision（视力）}
\ex{The glasses corrected her poor vision completely.}
\col{tunnel vision（隧道视野，喻只见一端）／ within one's vision（在视野内）}
\ex{The committee's tunnel vision missed the real problem.}
\col{a vision / a visionary（幻象／有远见者）}
\ex{In his vision, the desert city was green with gardens.}

\begin{gram}{vis-“看”词族收官（三访）}
\textbf{vision} $\leftarrow$ 拉丁语 \textit{vidēre}“\textbf{看}”——本讲义“看”家谱收官：希腊支 \textit{idein/eidos}（\textbf{idea} 理念，见 ideal 词条）、日耳曼支 \textit{wit/wise}（\textbf{wise} 明智、\textbf{witness} 见证人）、拉丁支 \textit{vidēre/vis-} 全家：\textbf{vision}（看见之力$\to$视力；远见）、\textbf{visual}（视觉的：\textit{visual arts}）、\textbf{evident}（e- 看得出$\to$明显的；\textbf{evidence}）、\textbf{provide}（预先看见$\to$提供；预防）、\textbf{supervise}（super- 从上看$\to$监督）、\textbf{revise}（re- 再看$\to$修订；复习）、\textbf{devise}（分开看$\to$设计）、\textbf{envisage / envision}（心中看见$\to$设想）、\textbf{interview}（inter- 互见$\to$面试；采访）、\textbf{viewpoint / worldview}（视点／世界观）。从“看”到“知”，三大文明用同一只眼睛。
\end{gram}

\vw{vitality}{/vaɪˈtæləti/}{n.}

\sense{n.}{生命力，活力，热情}
\col{be full of vitality（充满活力）}
\ex{Even in her eighties, Grandma is full of vitality.}
\col{economic / cultural vitality（经济／文化活力）}
\ex{Street markets give the city much of its economic vitality.}
\col{vital（形容词：生死攸关的；极其重要的；〔医〕生命体征的）}
\ex{Regular exercise is vital to a healthy heart.（be vital to）}
\ex{The nurse checked the patient's vital signs.（vital signs 生命体征）}

\begin{gram}{vit-“生命”词族}
拉丁语 \textit{vita}“\textbf{生命}”（原始印欧语 *\textit{gʷeih₃-}“活”，与英语 \textbf{quick} 的古义“活的”〔\textit{the quick and the dead} 生者与死者〕、\textbf{alive} 同根！）家族：\textbf{vital / vitality}（生命的／活力）、\textbf{vitamin}（\textit{vita} $+$ \textit{amine}胺——20 世纪波兰化学家 Funk 的杰作造词“生命之胺”）、\textbf{vitals}（要害；生命体征）、\textbf{revitalize}（re- 再赋生命$\to$使复兴：\textit{revitalize the old town}）——外加直接借入的拉丁短语 \textit{vita brevis, ars longa}（生命短暂，艺术长久）与 \textit{curriculum vitae}（“人生履历”$\to$简历，缩写 CV）。一词根，从药品柜到求职信。
\end{gram}

\vw{vivid}{/ˈvɪvɪd/}{adj.}

\sense{adj.}{鲜明的，耀眼的；生动的}
\col{a vivid memory / imagination（鲜活的记忆／想象）}
\ex{She has a vivid memory of her first day at school.}
\col{vivid colors / descriptions（鲜艳的色彩／生动的描写）}
\ex{The market was a blur of vivid colors and sharp smells.}
\ex{The novel's vivid descriptions of the famine shocked readers.}
\col{vividly（清晰地；生动地）}
\ex{I vividly remember the night the lights came back on.}

\begin{gram}{vivid / lively / lively-lookalikes 辨析与写作应用}
\textbf{vivid} 强调\textbf{如亲眼所见般的鲜明}（记忆、色彩、描写——读者“看得见”的文字）；\textbf{lively} 强调\textbf{活泼好动}（\textit{a lively child / discussion}）；\textbf{alive} 表语“活着的；热闹的”（\textit{the street came alive}）；\textbf{lifelike}“栩栩如生的”（艺术品逼真）。读后续写评分要点“描写生动”的地道实现即多用 vivid + 具体感官细节：\textit{a vivid red scarf, a vivid flash of lightning}。
\end{gram}

\begin{origin}
拉丁语 \textit{vividus}“\textbf{活的、有生气的}”，出自 \textit{vivere}“\textbf{活}”（与 vita“生命”同根，见上条；原始印欧语 *\textit{gʷeih₃-}）。\textit{vivere} 的家族在英语里个个精神：\textbf{vivid}（鲜活$\to$鲜明生动）、\textbf{survive}（sur- 超越而活$\to$\textbf{幸存}：\textit{survive the earthquake}；\textit{survival of the fittest} 适者生存）、\textbf{revive}（re- 复活$\to$苏醒；复兴：\textit{revive the economy}）、\textbf{vivacious}（活泼动人的，人如其词）、\textbf{vivisection}（活体解剖）、法语的 \textbf{bon vivant}（享乐生活家）、\textbf{viva}（口试，英式大学俚语“活考”）。“生动”的词源就是“\textbf{活着}”——文字活了，读者才看得见。
\end{origin}

\vw{volume}{/ˈvɒljuːm/}{n.}

\sense{n.}{音量，响度；体积，容积，容量；量，额；（书之）卷}
\col{turn the volume up / down（调高／调低音量）}
\ex{Please turn the volume down --- the neighbors are sleeping.}
\col{the volume of traffic / trade / sales（交通量／贸易额／销量）}
\ex{The new bridge doubled the volume of traffic across the river.}
\col{the volume of a container（容积，几何）}
\ex{Calculate the volume of the cylinder in cubic meters.}
\col{a volume of poetry（一卷诗集；volumes 大量）}
\ex{She spoke volumes about their friendship without saying a word.（见 speak volumes 词条）}

\begin{gram}{volume 的四义同源与“量”的选词}
四义（音量／体积／量额／卷册）皆出拉丁语 \textit{volumen}“卷起的书”（词源详见 speak volumes about 词条）：书卷$\to$份量$\to$体积$\to$音量，一条“由书到声”的引申链。科技写作的“量”各有专职：\textit{volume}（体积、总量）、\textit{amount}（数量，配不可数）、\textit{quantity}（量，数理）、\textit{capacity}（\textbf{容量}能力：\textit{the stadium's capacity} 体育场容纳量）。勿以 volume 泛代诸“量”——选词如选容器。
\end{gram}

%================================
\section{W}
%================================

\vw{wage}{/weɪdʒ/}{n.\ \&\ vt.}

\sense{n.}{（通常指按周领的）工资，工钱}
\col{a weekly / minimum / living wage（周薪／最低工资／生活工资）}
\ex{The city raised the minimum wage to 25 yuan an hour.}
\col{wages（复数泛指工资收入）／ wage-earner（养家糊口者）}
\ex{Both parents are wage-earners in most families here.}

\sense{vt.}{开始，发动（战争、运动）}
\col{wage war on / against（向……开战；大力开展）}
\ex{The government waged war on waste and corruption.}
\col{wage a campaign / battle}
\ex{Local groups waged a long campaign to save the wetlands.}

\begin{gram}{wage / salary / pay / income 工酬四分}
\textbf{wage}（蓝领\textbf{周薪／时薪}，恒多复数 wages）；\textbf{salary}（白领\textbf{月薪年薪}，固定总额）；\textbf{pay}（统称，不可数：\textit{equal pay for equal work} 同工同酬；\textit{pay rise} 加薪）；\textbf{income}（一切\textbf{收入}总称，含投资收益：\textit{income tax} 个人所得税）。搭配锁：\textit{draw a salary} 领薪、\textit{earn wages} 挣工钱、\textit{live on a fixed income} 靠固定收入生活——人物的阶级与职业，常藏在他领哪一种“钱”里。
\end{gram}

\begin{origin}
名词 \textbf{wage} 与动词 \textbf{wager}（打赌）同出古法语 \textit{gage / wagier}“\textbf{抵押、担保}”$\leftarrow$ 拉丁语语境的 *\textit{wadium}“\textbf{质押物}”（日耳曼词根 *\textit{wadja-}，经高卢拉丁语“倒灌”入罗曼语，再由诺曼法语带回英语）。“工资”的本义是\textbf{押给雇工的担保品}$\to$按约支付的报酬；“发动战争”（wage war）的本义是\textbf{把己方押上赌局}——战争是最豪险的 wager（\textbf{double or nothing}）。同根的还有 \textbf{engage}（en- 把质押投入$\to$\textbf{订约；订婚}；\textit{engaged} 忙着的、占线的）、\textbf{mortgage}（mort 死 + gage 押“死押”$\to$\textbf{按揭}——抵押到期不赎即“死”约）。一族的命运：钱、婚、房、战，全靠一个“押”字。
\end{origin}

\vw{wander}{/ˈwɒndər/}{vi. \& vt.}

\sense{vi. \& vt.}{闲逛，漫步；走失；走神；蜿蜒曲折}
\col{wander around / through（闲逛）}
\ex{We wandered through the old town for the whole afternoon.}
\col{wander off / away（走散、走失）}
\ex{The little boy wandered off while his mother was paying.}
\col{one's mind / thoughts wander（走神）}
\ex{Halfway through the lecture, his mind began to wander.}
\col{wander off the subject / point（离题）}
\ex{The essay wanders off the point in its final pages.}
\col{winding paths wander（路、河蜿蜒）}
\ex{A narrow stream wanders through the valley floor.}

\begin{gram}{wander / wonder：一对最容易互换的邻居}
\textbf{wander} /ˈwɒndər/“漫游、走神”（身体的迷路）；\textbf{wonder} /ˈwʌndər/“惊奇；想知道”（心智的迷失与追问）。两词在古英语中竟真是一家（\textit{wandrian} 与 \textit{wundor} 同根 *\textit{wen-}? 学界认为 wander$<$*\textit{wandrōjaną}“转向”、wonder$<$*\textit{wundrą}“奇物”，属\textbf{同源分家}），后分道扬镳——\textbf{身体漫步的走神}与\textbf{心灵惊奇的出神}本是同一种“出离”。拼写只差一字母 a/u，听辨靠元音；作文互错属低级错误，检查时专扫这一对。习语 \textit{wonder wander}：“且行且惊”是旅行文学的双关常客。
\end{gram}

\vw{weep}{/wiːp/}{vi. \& vt.\ (wept, wept)}

\sense{vi. \& vt.}{哭泣，流泪}
\col{weep at / over sth（因……而落泪）}
\ex{She wept openly at the news of the rescue.}
\ex{He wept over his lost youth at the reunion.}
\col{weep for / with joy（喜极而泣）}
\ex{The crowd wept with joy as the sailors came ashore.}
\col{weeping willow（垂柳——枝条如泪）；weep one's heart out（哭得死去活来）}
\ex{A weeping willow trailed its branches in the pond.}

\begin{gram}{weep / cry / sob / wept 型变化}
强度与方式：\textbf{weep}（无声或少声地\textbf{流泪}，书面、庄重）；\textbf{cry}（出声地哭，通用）；\textbf{sob}（抽泣、呜咽）；\textbf{wail}（嚎啕）；\textbf{burst into tears}（夺泪而出的一瞬）。变位 \textbf{weep--wept--wept} 与 sweep/sleep/keep/feel 同属“长元音 + t”强变化类（见 sweep 词条）。文学传统中 weep 自带悲剧腔——莎士比亚舞台指示 \textit{She weeps.} 与 \textit{She cries.} 的情感重量判然有别；读后续写哭戏按“轻 weep、中 cry、重 sob/wail”配药。
\end{gram}

\begin{origin}
古英语 \textit{wēpan}“哭泣、哀号”（原始日耳曼语 *\textit{wōpjanan}“喊叫”，与古诺斯语 \textit{óp}“呼喊”同根——\textbf{weep 的祖先竟是“喊”而不是“流泪}”：哭泣与呼号共享一个声学的词根，先民的悲声是不分泪与喊的）。它的表亲 \textbf{wap}（古语哀叹）已亡，而 \textbf{woe}（悲哀：\textit{woe is me}！古典哀叹）与之遥遥呼应（另源但同气）。对比拉丁语 \textit{plangere}（捶胸而哭$\to$\textbf{plangent} 凄厉作响的；\textbf{plaintiff} 原告“哀诉者”）与希腊语 \textit{klaiein}（$\to$\textbf{clamor}? 另源）——各文明的“哭”，各有各的声学。
\end{origin}

\vw{whisper}{/ˈwɪspər/}{vi.\&vt.\ \&\ n.}

\sense{vi. \& vt.}{耳语，低语；私下说；沙沙作响}
\col{whisper to sb（对某人耳语）／ whisper sth in sb's ear}
\ex{She leaned over and whispered the answer in my ear.}
\col{whisper about（私下议论，含传播义）}
\ex{Villagers whispered about the strange lights on the hill.}
\col{the wind / leaves whisper（风、叶沙沙作响——拟人）}
\ex{The bamboo leaves whispered in the night breeze.}

\sense{n.}{低语；轻柔的声音；传言}
\col{in a whisper（低声地）／ speak in whispers}
\ex{``He's gone,'' she said in a whisper.}
\col{a whisper of sth（一丝轻微声响／迹象）}
\ex{We heard a whisper of music from the next room.}
\col{whispers（流言）／ without a whisper of（毫无……之意）}
\ex{He agreed without a whisper of complaint.}

\begin{gram}{拟声词（onomatopoeia）小集}
\textbf{whisper} 属 \textit{wh-} 气音家族——wh- 开头的词多与\textbf{气息}相关：\textbf{whistle}（哨音）、\textbf{whisk}（拂扫）、\textbf{whirl}（回旋）、\textbf{whiff}（一缕气味）、\textbf{wheeze}（喘息）。广义拟声词是文学描写的音效库：\textit{the leaves whispered, the gate creaked, the thunder boomed}（回访 boom 词条）。读后续写的高级感一半来自“\textbf{以声写静}”：\textit{The silence was so complete that I could hear the candles whisper.}
\end{gram}

\vw{wind}{/waɪnd/}{vi.\&vt.\ (wound, wound)}

\sense{vi. \& vt.}{蜿蜒，迂回；上发条；缠绕}
\col{wind one's way（蜿蜒前行，叙景高频）}
\ex{The old silk road winds its way through deserts and mountains.}
\ex{A path wound down to the fishing village.}
\col{wind up（上紧发条；以……告终，口语）}
\ex{He forgot to wind up the clock, and it stopped at three.}
\ex{If you keep borrowing, you'll wind up in debt.（落得……下场）}
\col{wind sth around sth（缠绕）}
\ex{She wound the scarf twice around her neck.}

\begin{gram}{wind / wound / wound 与 wind / wound / wounded 的双胞胎陷阱}
动词 \textbf{wind}“蜿蜒、上发条”读 /waɪnd/，变位 wind--wound--wound（\textbf{wound} 读 /waʊnd/）；而名词 \textbf{wind}“风”读 /wɪnd/，另有动词 \textbf{wound}“使受伤”读 /wuːnd/，其过去式 \textbf{wounded}。于是 \textit{The road wound up the hill}（路蜿蜒上山，/waʊnd/）与 \textit{He was wounded}（负伤，/wuːnd/）字形近而音义皆异——朗读、听力的顶级陷阱。速判：\textbf{wound} 若讲“路、河、发条、缠绕”读 /waʊnd/，讲“受伤”读 /wuːnd/；\textit{wound up} 多半与发条或结局有关，与伤口无关。
\end{gram}

\begin{origin}
动词 \textbf{wind}“缠、弯”$\leftarrow$ 古英语 \textit{windan}“扭转、缠绕”（原始日耳曼语 *\textit{windaną}，与 \textbf{wander}“漫游”可能同源“转来转去”！见 wander 词条；德语 \textit{winden}）。家族：\textbf{wind up}（缠紧发条$\to$告终）、\textbf{rewind}（回卷——磁带时代的遗产词，数字界面仍在用）、\textbf{unwind}（解开$\to$\textbf{放松}， unwind after work）、\textbf{windmill}（风车“风之缠绕坊”？实为 wind 风 + mill 磨坊，\textbf{两个 wind 在此相遇}：风车既属 wind\textsubscript{1}（风）的产业，又靠 wind\textsubscript{2}（转动）工作）、\textbf{winding stairs}（旋梯）。名词 \textbf{wind}“风”（\textit{winds} 风向；\textit{windy} 多风的）另属拉丁 \textit{ventus} 家族的日耳曼表亲（\textbf{ventilate} 通风、\textbf{vent} 排气口与之遥望）。
\end{origin}

\vw{worthy}{/ˈwɜːði/}{adj.}

\sense{adj.}{值得（或应得）……的；值得尊敬的，值得注意的}
\col{be worthy of sth（固定搭配：值得……，接名词）}
\ex{Her devotion to the students is worthy of the highest praise.}
\ex{The proposal is worthy of serious consideration.}
\col{a worthy cause / opponent / successor（可敬的／旗鼓相当的）}
\ex{The runners-up proved worthy opponents to the champions.}
\col{worthy to do sth（配得上做）}
\ex{He proved himself worthy to wear the captain's armband.}
\col{worthlessness?（注意构词：worthless 一文不值的；unworthy 不配的）}
\ex{The fake painting turned out to be worthless.}

\begin{gram}{worth / worthy / worthwhile 三分}
（1）\textbf{worth}（介词性）：\textit{The museum is \textbf{worth} a visit / worth visiting.}（后接名词或动名词\textbf{主动式表被动}——与 deserve doing 同律，见 deserve 词条）；\ding{55} \textit{worth being visited}。（2）\textbf{worthy}（形容词）：须搭 \textit{of}，\textit{worthy of praise / worthy to be praised}（不定式则用被动）。（3）\textbf{worthwhile}（形容词）：“值得花时间做的”：\textit{a worthwhile project / It is worthwhile to compare the two designs.}三词搭配各异、互换必错，是语法填空的钉子户；一并记住 \textit{worth one's while}（不辜负某人所费）。
\end{gram}

\begin{origin}
\textbf{worthy} $=$ \textbf{worth}（价值）$+$ \textbf{-y}（充满……的）——古英语 \textit{weorþ}“价值、配得”（原始日耳曼语 *\textit{werþaz}，与原始印欧语 *\textit{wert-}“转”相关：价值即“\textbf{交换中转手之物}”——恰与 version 词条的“转”遥相握手！）。同根家族：\textbf{worship}（\textit{weorþscipe}“\textbf{worth-ship}价值之尊”$\to$\textbf{崇拜}——敬拜的本义是“承认对方的价值”！）、\textbf{worthwhile}（while“值得花的时间”）、\textbf{worthy / unworthy}、德语 \textbf{Wert}。一句词源，胜过十篇伦理学论文：所谓崇拜（worship），就是把价值（worth）还给应得者（the worthy）。
\end{origin}

\end{document}
